\chapter{Abstract}
\begin{SingleSpace}

\noindent \textbf{Title:} Development of a System for Analyzing Data on Network Attacks to Create Effective Strategies to Protect Against New Threats.\\
\textbf{Author:} Milana Dulatovna Muratbekova.\\
\textbf{Year:} 2026.

\vspace{0.8em}
\noindent Cyberattacks on public sector organizations, banks, and critical infrastructure have increased exponentially. Therefore, intelligent and adaptive intrusion detection systems are needed to detect attacks.
However, traditional signature and rule based security systems do not adequately detect new types of attacks. In addition, the use of artificial intelligence technology makes it difficult to develop new signature and rules. Data driven and machine learning based approaches are being developed to improve detection capabilities.

\vspace{0.8em}
\noindent Therefore, this thesis includes the design, implementation and testing of a complete machine learning pipeline to analyze network attack data.
A pre-processing step was used to clean up the cicids2017 dataset which is provided by the Canadian Institute for Cybersecurity. The pre-processing included removing duplicate records, standardizing scores using z-score normalization method, converting the highly skewed features into logarithms, selecting features using correlation and variance inflation factor (vif) techniques to reduce dimensionality of 80 to 74 features, and class imbalance correction using the smote technique combined with random undersampling.

\vspace{0.8em}
\noindent Six supervised classifiers were designed and tested: Logistic Regression, k-nearest neighbors (KNN), Random Forest, XGBoost, multilayer perceptron (mlp) classifier and lightgbm. The lightgbm gradient boosting model obtained the highest accuracy value across all measures of accuracy: accuracy = ~0.98; F1 score = ~0.97; and ROC-auc = ~0.99. Additionally, lightgbm had the lowest training time among the three gradient boosting models tested (lightgbm = ~89 seconds vs. XGBoost = ~317 seconds).
These results confirm the advantages of lightgbms' leaf-wise tree growing approach and histogram-based gradient calculations when applied to large dimensional table-based network traffic data.

\vspace{0.8em}
\noindent Shap global feature importance analysis and local interpretable model-agnostic explanations (lime) instance level explanations were used to evaluate model interpretability and transparency. four key discriminant flow features of the network were determined to be flow duration, incoming packet rate (inpktrate), Outgoing Byte Rate (outbyterate), and number of concurrent connections made by the same source ip address. These findings align closely with the existing body of literature within the field of cybersecurity and provide a justifiable foundation for the creation of alerts.

\vspace{0.8em}
\noindent Unsupervised anomaly detection using isolation forests and temporal modeling using long short-term memory (lstm) networks were used to identify previously unknown attack patterns and demonstrate their complementary role within hybrid detection architecture.

\vspace{0.8em}
\noindent A system architecture was proposed that integrates the trained model with commercial SIEM platforms to enable real-time explanation-enabled detection and alerting as well as integration with workflows related to incident management.
The study supports the central hypothesis: that applying machine learning to network traffic data analysis yields accurate, interpretable, and operationally deployable intrusion detection systems that generate viable defense plans for previously unknown or emerging cyber threats.

\vspace{0.8em}
\keywords{network intrusion detection, machine learning, LightGBM, XGBoost, SHAP, LIME, cybersecurity, anomaly detection, SIEM integration, gradient boosting, feature importance, network traffic analysis.}

\vspace{0.8em}
\noindent \textbf{Volume:} 5 chapters, 21 figures, 10 tables, 21 references, 52 pages.

\end{SingleSpace}
